%
% teil2.tex -- Beispiel-File für teil2 
%
% (c) 2020 Prof Dr Andreas Müller, Hochschule Rapperswil
%
% !TEX root = ../../buch.tex
% !TEX encoding = UTF-8
%
\section{Konvergiert das Basel-Problem überhaupt?}
\label{basel:section:konvergenz}
 
Um die Konvergenz bestimmen zu können, braucht man die Ungleichung
\[
    \frac{1}{n^2} < \frac{1}{n(n-1)},
\]
die für alle natürlichen Zahlen $n > 1$ gilt.
Daraus folgt die Ungleichung für Reihen
\[
    \sum_{n=2}^{\infty} \frac{1}{n^2}
    \leq
    \sum_{n=2}^{\infty} \frac{1}{n(n-1)}
    =
    \sum_{n=1}^{\infty} \frac{1}{n(n+1)},
\]
bei der man darauf achten muss, dass die Startwerte der Laufvariablen so gesetzt werden,
dass man nicht durch Null dividiert.
Dies spielt aber für die Konvergenzüberlegung keine Rolle.
 
Zur Berechnung schaut man sich die $m$-te Partialsumme an,
\[
    \sum_{n=1}^{m} \frac{1}{n(n+1)}
    =
    \frac{1}{1 \cdot 2} + \frac{1}{2 \cdot 3} + \frac{1}{3 \cdot 4} + \cdots + \frac{1}{m(m+1)},
\]
und schreibt jeden Summanden als Differenz zweier Brüche.
So erhält man die folgende \emph{Teleskopsumme}
\begin{align*}
    \left(1 - \frac{1}{2}\right)
    + \left(\frac{1}{2} - \frac{1}{3}\right)
    + \left(\frac{1}{3} - \frac{1}{4}\right)
    + \cdots
    + \left(\frac{1}{m} - \frac{1}{m+1}\right).
\end{align*}
Diese heisst so, weil sich zwei benachbarte Terme neutralisieren und man somit die Reihe zusammenfahren kann.
So bleibt nur noch der erste und der letzte Term
\[
    1 - \frac{1}{m+1}
\]
übrig.
Lässt man $m$ gegen unendlich gehen, geht diese Differenz gegen eins,
\[
    1 - \frac{1}{m+1} \xrightarrow{m \to \infty} 1.
\]
Anhand des sogenannten \emph{Majorantenkriteriums} weiss man also, dass die Reihe des Basel-Problems konvergiert.
Ebenfalls weiss man jetzt, dass die Summe der Reihe nicht grösser als zwei sein kann.